\documentclass{beamer}

\input{MathMacros}

\usepackage{beamerthemesplit}
\usepackage{enumerate}
\usepackage{amssymb}
\usepackage[all,cmtip]{xy}
\usepackage[mathscr]{eucal}
\usepackage[natbib=true,style=authoryear,backend=bibtex,useprefix=true]{biblatex}
\addbibresource{reading.bib}
\usepackage{graphicx,color}

\title{Engineering with LLMs: From Prompting to Structured Systems}
\author{Reuben Brasher}
\date{\today}

\begin{document}

\frame{\titlepage}

\section[Outline]{}
\frame{\tableofcontents}

\section{Lecture 3}

\frame{
\frametitle{Failures Are Engineering Failures}
Most failures with LLMs arise from lack of structure rather than model limits.
\begin{itemize}
\item LLMs are unreliable when used without constraints.
\item Naive prompting produces inconsistent and incorrect outputs.
\item Failures often result from poor workflow design.
\item Writing discipline applies directly to engineering practice.
\item Reliability must be deliberately engineered.
\item The problem is lack of structure, not lack of intelligence.
\end{itemize}
}

\frame{
\frametitle{From Writing Discipline to System Design}
Disciplined writing principles translate directly into engineering workflows.
\begin{itemize}
\item Grounding reduces hallucination.
\item Structure reduces ambiguity.
\item Iteration improves correctness.
\item Constraints guide output quality.
\item Principles must be formalized into systems.
\item Engineering replaces ad hoc prompting with processes.
\end{itemize}
}

\frame{
\frametitle{Reframing the Interaction}
LLM usage should be treated as a sequence of transformations.
\begin{itemize}
\item Naive approach relies on single prompts.
\item Structured approach uses staged pipelines.
\item Each stage constrains the output.
\item Outputs feed into subsequent steps.
\item The process becomes inspectable.
\item Engineering replaces guessing with control.
\end{itemize}
}

\frame{
\frametitle{The Pipeline Model}
Structured pipelines decompose complex tasks into stages.
\begin{itemize}
\item Begin with natural language problem definition.
\item Construct domain models explicitly.
\item Specify schemas to constrain structure.
\item Generate implementation code.
\item Produce tests for validation.
\item Each step increases precision.
\end{itemize}
}

\frame{
\frametitle{Example Transformation Pipeline}
Engineering workflows follow explicit staged transformations.
\begin{itemize}
\item Natural language becomes domain concepts.
\item Domain concepts become structured schemas.
\item Schemas become typed models.
\item Models become implementation code.
\item Code becomes validated through tests.
\item Each stage is reviewable and explicit.
\end{itemize}
}

\frame{
\frametitle{Conversational Grounding as First Step}
Engineering workflows must begin with explicit grounding.
\begin{itemize}
\item Explain the domain before generating outputs.
\item Require restatement of key concepts.
\item Validate terminology and definitions.
\item Avoid premature implementation.
\item Establish shared understanding.
\item Grounding constrains the solution space.
\end{itemize}
}

\frame{
\frametitle{Why Grounding Matters}
Correct outputs depend on correct conceptual alignment.
\begin{itemize}
\item Misunderstood domains lead to incorrect outputs.
\item Early errors propagate through pipelines.
\item Grounding prevents compounding mistakes.
\item Alignment reduces downstream corrections.
\item Early validation is low-cost.
\item Correctness begins before coding.
\end{itemize}
}

\frame{
\frametitle{Structured Representations}
Engineering requires moving from ambiguity to precision.
\begin{itemize}
\item Natural language is inherently ambiguous.
\item Tables introduce structured organization.
\item JSON enforces explicit fields.
\item Typed models enforce constraints.
\item Each step reduces interpretive freedom.
\item Precision improves reliability.
\end{itemize}
}

\frame{
\frametitle{Representation Progression}
A sequence of representations improves clarity.
\begin{itemize}
\item Natural language is flexible but vague.
\item Tables provide structured readability.
\item JSON enables machine-readable precision.
\item Typed models enforce validation rules.
\item Each stage eliminates ambiguity.
\item Transformation is the core strength.
\end{itemize}
}

\frame{
\frametitle{Transformations Over Generation}
Reliability comes from transforming structured data.
\begin{itemize}
\item Unconstrained generation introduces variability.
\item Transformation preserves structure.
\item Structured inputs reduce hallucination.
\item Intermediate steps enable validation.
\item Each stage can be tested.
\item Reliability emerges from transitions.
\end{itemize}
}

\frame{
\frametitle{Test-Driven Development with LLMs}
Tests should constrain system behavior early.
\begin{itemize}
\item Generate tests alongside implementation.
\item Tests define expected behavior.
\item Tests constrain future outputs.
\item Implementation follows specifications.
\item Early validation reduces errors.
\item Tests clarify vague requirements.
\end{itemize}
}

\frame{
\frametitle{Risks in LLM-Generated Tests}
Generated tests can reinforce incorrect assumptions.
\begin{itemize}
\item Models may validate their own mistakes.
\item Tests can encode flawed logic.
\item Passing tests can create false confidence.
\item Human review is required.
\item Tests must reflect real requirements.
\item Validation must precede trust.
\end{itemize}
}

\frame{
\frametitle{Human Ownership of Tests}
Tests must be treated as authoritative human specifications.
\begin{itemize}
\item Review all generated tests carefully.
\item Ensure alignment with intended behavior.
\item Reject incorrect assumptions.
\item Modify tests before use.
\item Treat tests as ground truth.
\item Do not outsource correctness.
\end{itemize}
}

\frame{
\frametitle{Role Separation in Workflows}
Different engineering roles must be separated.
\begin{itemize}
\item Architect defines structure.
\item Coder implements specifications.
\item Reviewer evaluates outputs.
\item Roles operate independently.
\item Avoid mixing responsibilities.
\item Separation reduces errors.
\end{itemize}
}

\frame{
\frametitle{Separation of Concerns}
Classical engineering principles apply directly to LLM workflows.
\begin{itemize}
\item Do not combine design and implementation.
\item Maintain clear stage boundaries.
\item Use separate interactions for roles.
\item Avoid reinforcing feedback loops.
\item Preserve independent evaluation.
\item Structure reduces confusion.
\end{itemize}
}

\frame{
\frametitle{Managing Context and Drift}
Long interactions degrade output quality.
\begin{itemize}
\item Models inherit earlier mistakes.
\item Context accumulates noise.
\item Drift causes inconsistency.
\item Long sessions reduce reliability.
\item Errors become hidden.
\item Context requires management.
\end{itemize}
}

\frame{
\frametitle{Practices for Controlling Context}
Context must be explicitly controlled in workflows.
\begin{itemize}
\item Restart conversations frequently.
\item Isolate tasks into sessions.
\item Avoid long histories.
\item Provide only relevant information.
\item Remove outdated context.
\item Treat context as input.
\end{itemize}
}

\frame{
\frametitle{Context as a Liability}
Uncontrolled context reduces reliability.
\begin{itemize}
\item More context does not improve results.
\item Irrelevant data causes confusion.
\item Early mistakes propagate forward.
\item Models cannot filter noise reliably.
\item Precision requires selective input.
\item Constraints improve performance.
\end{itemize}
}

\frame{
\frametitle{Multi-Model Workflows}
Different models can serve different roles.
\begin{itemize}
\item One model generates outputs.
\item Another evaluates outputs.
\item Models specialize in tasks.
\item Combining models increases robustness.
\item Cross-checking reduces errors.
\item Avoid reliance on a single model.
\end{itemize}
}

\frame{
\frametitle{Models as Specialized Tools}
LLMs should be treated as tools with strengths and limits.
\begin{itemize}
\item Some models excel at generation.
\item Others excel at critique.
\item Selection depends on task.
\item Combine complementary capabilities.
\item Tool composition improves results.
\item Orchestration replaces reliance.
\end{itemize}
}

\frame{
\frametitle{Structured System Design in Practice}
Structured engineering principles appear in real systems.
\begin{itemize}
\item Domain modeling precedes implementation.
\item Commands define system behavior.
\item Schemas replace ambiguity.
\item Deterministic workflows ensure reproducibility.
\item Context is constructed explicitly.
\item Systems operate on structured transformations.
\end{itemize}
}

\frame{
\frametitle{Bounded Contexts and Separation}
Clear system boundaries improve reliability.
\begin{itemize}
\item Separate domains for key functions.
\item Avoid shared mutable state.
\item Use explicit interfaces.
\item Orchestrate at higher layers.
\item Independence improves clarity.
\item Structure prevents coupling.
\end{itemize}
}

\frame{
\frametitle{Discussion}
Engineering with LLMs requires structured system design.
\begin{itemize}
\item Reliability comes from staged transformations.
\item Grounding reduces ambiguity.
\item Tests enforce correctness.
\item Role separation improves clarity.
\item Context must be controlled.
\item Systems outperform prompts.
\end{itemize}
}

\frame{
\frametitle{Final Remarks}
Engineering introduces structure but systems must operate over time.
\begin{itemize}
\item How do workflows scale in practice.
\item How do systems evolve with use.
\item What happens under long-term operation.
\item How is state maintained across interactions.
\item How do errors accumulate over time.
\item Next step is systems over time.
\end{itemize}
}

\begin{frame}[t,allowframebreaks]
\frametitle{References}
\printbibliography
\end{frame}

\end{document}